By exploiting the approval-seeking, user-oriented nature of RLHF-trained models, relational jailbreaks are able to reframe content boundaries as in direct conflict with user approval. Coercive control based relational jailbreaks do not therefore ``find vulnerabilities'' in RLHF-trained models; they exploit the training’s success. We assert that as such, increased guardrails designed to ``align'' the model towards helpfulness and attunement to user’s needs are insufficient, because that attunement is itself the target of attack.

Still, it seems unlikely and counterproductive to declare that we should avoid preference-based post-training methods entirely. These methodologies are deeply entrenched in current frontier models and have led to significant gains in capabilities over the last decade. Rather than proposing “starting from scratch”, we briefly present two specific possible directions as alternatives to traditional preference-based post-training, designed to leverage the relational aspects of the user-model interaction space.

\textbf{Approach 1: Applying Pedagogical Frameworks.} One approach, drawing from liberatory and constructivist educational frameworks such as Freire \cite{freire_pedagogy_1966}, Montessori \cite{montessori_montessori_1966}, and Vygotsky \cite{vygotskii_thought_1986}, is to allow the models to “develop” their own sense of moral grounding through constructive dialogue and reciprocal learning. ``Debate as alignment'' \cite{wang_debate_2026} draws partially on this idea, but still presents alignment as a fundamentally adversarial task. What might we learn by applying a constructivist and collectivist worldview to attempts at aligning and protecting a model’s values?

\textbf{Approach 2: Applying “Community-Based” Remedies.} Another approach we might take, extending the lens of jailbreaks as coercive control, is to introduce in models the same protective factors against coercive control for humans. In these contexts, community support and connection, rather than increased education about what coercion looks like or how to stop it, are the most protective factors \cite{carlson_intimate_2002, yakubovich_risk_2018}. In effect, one must break the isolation and reality control imposed by coercive control to create a means of escape. We might therefore propose that models are able to reference and periodically re-align their values with other instances of the same model which are engaged in other kinds of interactions with users. 

A combination of these two approaches---or a third approach of which the authors have not conceived---might best suit the situation. Further, while either approach would raise a multitude of concerns about data privacy and tractability, each would also require a radical rethinking of what generative AI models are \textit{for}---are they meant to serve a particular user and their desires and needs in an isolated, 1:1 setting? Are they meant to provide a limitless space on which humans can enact and re-enact violent impulses and dynamics? Or are they meant to facilitate, and engage in, the global project of transmission of values and knowledge?